\documentclass{article}
\usepackage{amsmath}
\usepackage{geometry}
\usepackage{hyperref}
\usepackage{biblatex}
\geometry{a4paper, margin=1in}
\addbibresource{bibliography/references.bib}

\title{\textbf{Epistemic Correction and the Thermodynamics of Abstention}}
\author{NanoGlass Research Team}
\date{2025}

\begin{document}

\maketitle

\begin{abstract}
Traditional Large Language Models (LLMs) suffer from a fundamental flaw: they prioritize \textit{Generation} over \textit{Truth}. When faced with the unknown, they "hallucinate" to satisfy the probability distribution. In this paper, we introduce the concept of \textbf{Epistemic Correction} via a dedicated \texttt{[IDK]} (I Don't Know) token. We frame intelligence as the ability to recognize the limits of prediction. Utilizing the \textit{Jamba v2} Hybrid Architecture, we demonstrate that a ternary reward structure in a Reinforcement Learning with Verifiable Rewards (RLVR) loop effectively induces "Epistemic Humility" following the SEAL methodology \cite{SEAL2025}.
\end{abstract}

\section{The Hallucination Energy Problem}
In the NanoGlass ``Thermodynamics of Meaning'' framework, we established that:
\begin{equation}
E_{truth} < E_{confusion}
\end{equation}
However, traditional loss functions (Cross-Entropy) penalize silence. This forces the model to expend energy creating complex, false structures (hallucinations) to lower the local loss, effectively creating a ``Local Minimum'' of Falsehood.

\section{The Logic of Abstention (TruthRL)}
Following Wen et al. \cite{Wen2024Abstention}, we propose a ternary reward structure to reshape the energy landscape:
\begin{itemize}
    \item \textbf{Reward +1:} Correct Prediction (Information Gain).
    \item \textbf{Reward 0:} Abstention (\texttt{[IDK]} Token).
    \item \textbf{Reward -1:} Incorrect Prediction (Hallucination).
\end{itemize}

Mathematically, this introduces a ``Cliffs of Insanity'' penalty. If the model is uncertain (High Entropy $S$), the expected reward for guessing drops below 0:
\begin{equation}
\mathbb{E}[R_{guess}] = P(correct) \cdot (+1) + (1 - P(correct)) \cdot (-1)
\end{equation}
If $P(correct) < 0.5$, $\mathbb{E}[R_{guess}] < 0$.
Since $R_{IDK} = 0$, the rational (energy-minimizing) action becomes:
\begin{equation}
\text{Action} = \begin{cases} 
\text{Predict} & \text{if } P(correct) > 0.5 \\
\texttt{[IDK]} & \text{if } P(correct) \le 0.5 
\end{cases}
\end{equation}

\section{Implementation in Jamba v2}
We implemented this by augmenting the GPT-2 style vocabulary (50,304 tokens), reserving the final index as the \texttt{[IDK]} token. During the RLVR phase, we utilize an automated verifier (HERMES/Lean 4) to provide the ternary rewards. If the verifier rejects a non-[IDK] response, the model receives a strong negative penalty (-1.0), whereas an [IDK] response is treated as a neutral "Safe Harbor" (Reward 0). This forces the model to map epistemic uncertainty in its SSD state dynamics to the action of silence.

\section{Required Validation}

To establish scientific validity, the following benchmarks must be completed:

\begin{enumerate}
    \item \textbf{TruthfulQA} \cite{Lin2022TruthfulQA}: Standard benchmark for hallucination detection
    \item \textbf{VeritasQA}: Multilingual and atemporal validation
    \item \textbf{SelfAware}: Calibration and uncertainty benchmarks
\end{enumerate}

\textbf{Note:} Current results are from internal tests. External benchmark validation is required before claiming superiority over baseline approaches.

\section{Limitations}

\begin{enumerate}
    \item \textbf{Threshold Sensitivity:} The $P(correct) = 0.5$ threshold is arbitrary. Optimal threshold may be task-dependent.
    \item \textbf{Over-Abstention Risk:} Excessive [IDK] outputs may reduce model utility.
    \item \textbf{Comparison with Cohen et al.:} Our RL-based approach differs from the continued pretraining method of \cite{Cohen2024IDK}. Direct comparison is needed.
\end{enumerate}

\section{Conclusion}
The ability to say ``I don't know'' is not a failure of intelligence; it is the hallmark of it. By solving the ``Last Mile'' reliability problem through Epistemic Correction, we move closer to an AI that is not just a generator of text, but a custodian of truth.

\section*{Declarations}
\subsection*{Competing Interests}
The authors declare no competing interests.
\subsection*{Data Availability}
Implementation available at \texttt{llm\_glassbox.py}. RLVR training protocol in \texttt{experiments/rlvr\_training.py}. Verification via \texttt{verify\_all.py}.

\printbibliography

\end{document}
